% !TeX root = main.tex

\chapter{Methodology}

\section{Language Choice}

The algorithms are implemented and tested in C\texttt{++} for several reasons.

\begin{itemize}
  \item A low-level compiled language gives you realistic estimates on how fast these algorithms
  actually are when they aren't held back by overheads like an interpreter.
  \item You get precise control over memory layouts and whether things go on the stack or the
  heap. For the most part, data structures like \texttt{std::array} are preferred because
  they reduce pointer chasing and are more cache-friendly.
  \item Multithreading is easy to perform in C\texttt{++}, thanks to tools like \texttt{std::jthread}.
  Although theoretical time complexities are important, whether or not an algorithm is practical
  and useful in real life also depends on how well it may be \textit{parallelised}.
\end{itemize}

\section{Point Generation}

A meaningful dataset should try and avoid degenerate cases like duplicate points. This
necessarily forces a large-enough range to circumvent the Birthday Paradox. In addition,
it should be easy to compute the distance between any two points, and compare distances
accurately. In other words, let's avoid floating point operations. Not only is it easy to mess up
comparisons, when adding floats together the ideal strategy is to sort them by size and add
the smallest numbers first before merging with the large ones. And this adds extra overhead
that detracts from the core task at hand.

So \texttt{int} it is, and because we're working with \textit{Euclidean} distances, a trick
to preserve precision and reduce operations is to avoid the square root altogether, and
compare squared distances instead. To spawn a set of points, we used the \textit{Mersenne Twister}
pseudorandom number generator (\code{std::mt19937_64}) to achieve a uniform random distribution
of points. Fixed seeds are also used to ensure fair benchmarking.
